% !PRINTED 390
\weberpage{390}
%<<CONT
einer Zahl in $A_1$ mit einer Zahl in $A_2$ darstellen, was wir symbolisch durch

\weq{2}{Z = A_1 A_2}

ausdrücken.

Der \gesperrt{Durchschnitt} $A$ der Gruppen $A_1$ und $A_2$ enthält alle und nur die Zahlen, die zugleich Normenreste von $m_1$ und von $m_2$ sind, und wir beweisen, daß $A$ die Gruppe der Normenreste $m$

\weq{3}{m = m_1 m_2}

ist.

Ist nämlich $a$ eine Zahl, die zugleich in $A_1$ und in $A_2$ enthalten ist, so gibt es zwei Zahlen $\omega_1, \omega_2$ in $O$, die den Bedingungen

\weq{4}{\begin{aligned} N(\omega_1) &\equiv a \pmod{m_1},\\ N(\omega_2) &\equiv a \pmod{m_2} \end{aligned}}

genügen. Es lassen sich nun zwei rationale Zahlen $x_1, x_2$ so bestimmen, daß

\[
\begin{aligned}
x_1 &\equiv 1 \pmod{m_1}, & x_2 &\equiv 0 \pmod{m_1},\\
&\equiv 0 \pmod{m_2}, & &\equiv 1 \pmod{m_2},
\end{aligned}
\]

und wenn wir dann also

\[ \omega = x_1 \omega_1 + x_2 \omega_2 \]

setzen, so ist

\weq{5}{\begin{aligned} N(\omega) &\equiv N(\omega_1) \equiv a \pmod{m_1},\\ &\equiv N(\omega_2) \equiv a \pmod{m_2}, \end{aligned}}

und da $m_1$ und $m_2$ relativ prim sind, so ist auch

\weq{6}{N(\omega) \equiv a \pmod{m},}

also $a$ Normenrest von $m$. Umgekehrt ist ein Normenrest von $m$ zugleich Normenrest jedes Teilers von $m$, also auch von $m_1$ und $m_2$.

Hieraus ergibt sich nun mit Hilfe der Sätze 13. und 14., \S\,100:

\[
\begin{aligned}
(Z, A_1) &= (A_1 A_2, A_1) = (A_1 A_2, A_1 A) = (A_2, A),\\
(Z, A) &= (Z, A_2)(A_2, A) = (Z, A_1)(Z, A_2),
\end{aligned}
\]

und wir haben die Beziehung:

\weq{7}{(Z, A) = (Z, A_1)(Z, A_2).}

Damit ist die Frage nach den Normenresten auf den Fall zurückgeführt, daß der Modul eine \gesperrt{Primzahlpotenz} ist.

\S\,108. \gesperrt{Normenreste der Primzahlpotenzen.}

Es sei $O$ eine Ordnungsgruppe im Körper $\Omega$ (\S\,98) und $P$, die Gruppen der rationalen Zahlen, die relativ prim zur Diskri-
%CONT>>
